\documentclass[10pt,a4paper]{amsart}
\usepackage[margin=19mm]{geometry}
\usepackage{amsmath,amssymb,mathtools}
\usepackage[scr=rsfs,cal=euler]{mathalfa}
\usepackage{xfrac}
\usepackage[unicode,hidelinks,bookmarks=false]{hyperref}
\newcommand{\ie}{i.e.\ }
\newcommand{\cf}{cf.\ }
\newcommand{\resp}{respectively\ }
\newcommand{\Subsection}{\subsection}
\providecommand{\nobreakdash}{\nobreak\hbox{-}}
\setlength{\emergencystretch}{2em}
\multlinegap0pt
\usepackage{bm}
\usepackage{bbm}
\usepackage{dsfont}
\usepackage{xspace}
\usepackage{cancel}
\usepackage{mathtools}
\usepackage[shortlabels]{enumitem}
\usepackage{amsthm}
\makeatletter
\@ifundefined{thm}{\newtheorem{thm}{Theorem}}{}
\makeatother
\newcommand{\bb}[1]{\mathbb{#1}}
\title[Existence and uniqueness of Remotely Almost Periodic solutio]{Existence and uniqueness of Remotely Almost Periodic solutions of differential equations with piecewise constant argument}
\author{Diego Jaure, Christopher Maulen}
\date{Source: arXiv:2602.00926v1 (CC BY 4.0)}
\begin{document}
\maketitle

\noindent\emph{Source and licence.} Existence and uniqueness of Remotely Almost Periodic solutions of differential equations with piecewise constant argument. Version arXiv:2602.00926v1 (CC BY 4.0), distributed under the Creative Commons Attribution 4.0 International licence, as recorded in the arXiv licence metadata for this exact version and shown on the abstract page https://arxiv.org/abs/2602.00926v1. Full rights notice: licenses/arxiv-2602.00926-CC-BY-4.0.txt.\par\smallskip

\noindent\textit{Editorial scope.} Counted unit: the Thm whose source label is \texttt{perturbed\_f\_nolineal} in the article (source0.tex lines 1312--1322, source label \texttt{perturbed\_f\_nolineal}), delivered together with the article's own complete original proof of it (source0.tex lines 1324--1400, one \texttt{proof} environment of 77 lines). No mathematics is written, repaired or re-derived: statement and proof are the article's own text line for line. Required context retained uncounted: DEPCA\_lineal (lines 119--123); DEPCA\_lineal+f (lines 119--123); intro\_perturbed (lines 170--172); teo\_perturbada\_depca (lines 175--186); H1 (lines 1021--1038); H2 (lines 1021--1038); eq:generadorF (lines 1295--1298); eq:perturbadof (lines 1295--1298); C1 (lines 1301--1310); C2 (lines 1301--1310); C3 (lines 1301--1310). Every definition, display and label of those lines is retained unchanged. Omitted from this package and not redistributed: 1--118, 124--169, 173--174, 187--1020, 1039--1294, 1299--1300, 1311--1311, 1323--1323, 1401--1523. Nothing inside a retained excerpt is omitted or altered: every retained body is delivered exactly as the article's lines are written, so no omission receipt of any kind accompanies this package. Citations: the retained excerpts carry no citation command at all, so no bibliography entry of the article is transcribed and no reference is redistributed. Reference adaptation: none was needed and none was made; every reference inside the delivered unit resolves inside it, so no unresolved reference or citation is produced. Printed slips kept literal: no printed word, symbol, hypothesis or number of the counted statement or of its proof was altered. Layout and class: the article's own class is \texttt{amsart} with packages including babel, fontenc, geometry, verbatim, amsthm, amssymb, setspace, enumerate, fancyhdr, enumitem, xcolor, csquotes, hyperref, biblatex; the wrapper is the lane's pinned \texttt{amsart} editorial wrapper at 10pt on A4, which restates the theorem-like environments the retained text needs and adds the article's own macro definitions that the retained text needs (\texttt{\textbackslash bb}), with extra wrapper packages (bm, bbm, dsfont, xspace, cancel, mathtools, enumitem). The delivered numbering is the unit's own. Assets: no retained span contains a figure environment, a graphics-inclusion command, a PDF-page inclusion, a table or a TikZ picture; no image file is copied into this package, the package declares no asset dependency and no image file is redistributed with it. Licence: version arXiv:2602.00926v1 of this article is distributed under the Creative Commons Attribution 4.0 International licence, as recorded in the arXiv licence metadata for this exact version and shown on the abstract page https://arxiv.org/abs/2602.00926v1. A full rights notice is delivered at licenses/arxiv-2602.00926-CC-BY-4.0.txt. Characters: every retained excerpt is pure ASCII, so the delivered engine, XeLaTeX with its default Latin Modern text font, reports no missing character.
\begin{eqnarray}
z'(t)&=& A(t)z(t) \label{DEPCA_A}\\
y'(t)&=& A(t)y(t)+B(t)y([t]) \label{DEPCA_lineal}\\
x'(t)&=& A(t)x(t)+B(t)x([t])+f(t), \label{DEPCA_lineal+f}
\end{eqnarray}

\begin{eqnarray}
y'(t) &=& A(t)y(t) + B(t)y([t]) + f(t) + g_{\nu}(t, y(t), y([t])), \label{intro_perturbed}
\end{eqnarray}

\begin{thm}\label{teo_perturbada_depca}
Let $\xi$ be the unique remotely almost periodic solution of \eqref{DEPCA_lineal+f}.  
If conditions \ref{H1} and \ref{H2} are satisfied,  
then there exists $r > 0$ and $\nu_0 > 0$ sufficiently small such that \eqref{intro_perturbed}  
has a unique remotely almost periodic solution $\psi_\nu$ in an $r$-neighborhood of $\xi$  
for each $\nu \in [0, \nu_0)$.  
Moreover, if $g_\nu(t, x, y)$ is uniformly continuous in $(t, x, y) \in \mathbb{R} \times B_r(0) \times B_r(0)$,  
 $\psi_\nu$ depends continuously on $\nu$ and  
\[
\lim_{\nu \to 0} \psi_\nu(t) = \xi(t), \quad \mbox{for all}~{} t\in\mathbb{R}.
\]
\end{thm}

\begin{enumerate}[label=\textbf{(H.\arabic*)}]
\item\label{H1} Let $A(t)$, $B(t)$, and $f(t)$ be remotely almost periodic functions.  
The discrete equation associated with \eqref{DEPCA_lineal} admits an exponential dichotomy such that its associated Green's kernel is bi-remotely almost periodic and summable.

\item\label{H2} Let $g_{\nu}(t, y_1, y_2)$ be a remotely almost periodic function in $t$  
uniformly for $(y_1, y_2, \nu) \in B_r(0) \times B_r(0) \times [0, \nu_0]$,  
and for each fixed small real parameter $\nu$, it is uniformly bounded.  
It satisfies the local Lipschitz condition  
\[
\left\Vert g_{\nu}(t, x_1, y_1) - g_{\nu}(t, x_2, y_2) \right\Vert \leq M_0(r, \nu) \left[ \left\Vert x_1 - x_2 \right\Vert + \left\Vert y_1 - y_2 \right\Vert \right],
\]
where $(t, x_1, y_1), (t, x_2, y_2) \in \mathbb{R} \times B_r(0) \times B_r(0)$ and $\nu \in [0, \nu_0]$,  
such that $M_0(r, \nu) \to 0$ and  
\[
\Vert g_{\nu} \Vert_r = \sup_{t \in \mathbb{R},\, \Vert x \Vert, \Vert y \Vert \leq r} \Vert g_\nu(t, x, y) \Vert \to 0
\]  
as $\nu \to 0$, for every fixed $r$.
\end{enumerate}

\begin{align}
\frac{dx}{dt} & =f\left(t,x(t),x([t])\right)\label{eq:generadorF}\\
\frac{dy}{dt} & =f\left(t,y(t),y([t])\right)+g_{\nu}\left(t,y\left(t\right),y([t])\right),\label{eq:perturbadof}
\end{align}

\begin{enumerate}[label=\textbf{(C.\arabic*)}]
\item\label{C1} The equation \eqref{eq:generadorF} has a unique remotely almost periodic solution $\xi(t)$.
\item\label{C2} $f(t, x, y) \in C^{2}$ in $x$ and $y$, and the second-order partial derivatives satisfy a Lipschitz condition in $x$ and $y$.
\item\label{C3} The discrete equation associated with the variational equation with piecewise constant argument
\begin{eqnarray*}
\frac{dz}{dt} = \frac{\partial f}{\partial x}(t, \xi(t), \xi([t])) z(t) + \frac{\partial f}{\partial y}(t, \xi(t), \xi([t])) z([t])
\end{eqnarray*}
admits an exponential dichotomy such that its associated Green's kernel is bi-remotely almost periodic and summable, where $\xi(t)$ is defined in \ref{C1}.  
Moreover, $\frac{\partial f}{\partial x}(t, x, y)$ and $\frac{\partial f}{\partial y}(t, x, y)$ are remotely almost periodic functions in the first variable.
\end{enumerate}

\begin{thm}\label{perturbed_f_nolineal}
If conditions \ref{C1}, \ref{C2} and \ref{C3}, along with \ref{H2} are satisfied,  
then there exist $r > 0$ and $\nu_0 = \nu_0(r)$ sufficiently small such that the system \eqref{eq:perturbadof}  
has a unique remotely almost periodic solution $\psi_\nu(t)$ in an $r$-neighborhood of $\xi(t)$, for all $\nu \in [0, \nu_0]$.

Moreover, if $g_\nu(t, x, z)$ is uniformly continuous on $(t, x, z) \in \mathbb{R} \times B[0, r] \times B[0, r]$, with $\nu \in [0, \nu_0]$,  
then $\psi_\nu(t)$ is continuous in $\nu$, and we have  
\[
\lim_{\nu \rightarrow 0} \psi_\nu(t) = \xi(t).
\]
\end{thm}

\begin{proof}
By differentiation, we know
\begin{equation}\label{calculo_aprox_f}
\begin{aligned}
f(t,x(t) &+\xi(t),y([t])+\xi([t]))-f(t,\xi(t),\xi([t])) 
\\
=&~{}\frac{\partial f}{\partial x}(t,\xi(t),\xi([t]))x(t) 
+\frac{\partial f}{\partial y}(t,\xi(t),\xi([t]))y([t]) \\
&+\frac{1}{2}(x(t)\cdot \nabla_{x}+y([t])\cdot \nabla_{y})^{2}f(t,\theta x(t)+\xi(t),\theta y([t])+\xi([t])) \ \ \ 
,
\end{aligned}
\end{equation}
where $\theta \in [0, 1)$, ``$\cdot$'' denotes the inner product, and ``$\nabla$'' is the Hamiltonian operator. \\  
Let $z(t) = y(t) - \xi(t)$. Then, from \eqref{eq:perturbadof} and \eqref{eq:generadorF}, we have that
\[
\begin{aligned}
\frac{dz}{dt}(t)=&~{}f(t,z(t)+\xi(t),z([t])+\xi([t]))-f(t,\xi(t),\xi([t]))\\
 & +\nu g\left(t,z(t)+\xi(t),z([t])+\xi([t])\right),
\end{aligned}
\]
by \eqref{calculo_aprox_f}, we have
\[
\begin{aligned}
\frac{dz}{dt}(t)=&\frac{\partial f}{\partial x}(t,\xi(t),\xi([t]))z(t)+\frac{\partial f}{\partial y}(t,\xi(t),\xi([t]))z([t])\\
 &+\frac{1}{2}(x(t)\cdot \nabla_{x}+y([t])\cdot \nabla_{y})^{2}f(t,\theta z(t)+\xi(t),\theta z([t])+\xi([t]))\\
 &+\nu g\left(t,z(t)+\xi(t),z([t])+\xi([t])\right).
\end{aligned}
\]
Let $A(t)=\frac{\partial f}{\partial x}(t,\xi(t),\xi([t])) $, $B(t)=\frac{\partial f}{\partial y}(t,\xi(t),\xi([t]))$ and 
\[
\begin{aligned}
F(t,z(t),z([t]))=&\frac{1}{2}(x(t)\cdot \nabla_{x}+y([t])\cdot \nabla_{y})^{2}f(t,\theta z(t)+\xi(t),\theta z([t])+\xi([t])),
\end{aligned}
\]
without loss of generality, we can assume that  
$\max \{ \Vert A \Vert_{\infty}, \Vert B \Vert_{\infty} \} \leq M$.  
We obtain the following system:
\[
\begin{aligned}
\frac{dz}{dt}(t)=A(t)z(t)+B(t)z([t])+F(t,z(t),z([t]))+\nu g\left(t,z(t)+\xi(t),z([t])+\xi([t])\right)
\end{aligned}
\]
By $(C_2)$, there are $r>0$ y $\nu^{0}$, $N_i(r)$, $i=1,2$ y $M_0(\nu_{0})$ such that
\[
\begin{aligned}
\left\Vert\frac{\partial^{2}f}{\partial x_{i}\partial x_{j}}(t,\theta z(t)+\xi(t),\theta z([t])+\xi([t]))\right\Vert\leq &~{} N_{1}(r)\\
\left\Vert\frac{\partial^{2}f}{\partial y_{i}\partial y_{j}}(t,\theta z(t)+\xi(t),\theta z([t])+\xi([t]))\right\Vert\leq &~{} N_{1}(r)\\
\left\Vert\frac{\partial^{2}f}{\partial x_{i}\partial y_{j}}(t,\theta z(t)+\xi(t),\theta z([t])+\xi([t]))\right\Vert\leq &~{} N_{1}(r), 
\end{aligned}
\]
for $i,j=1,2,\cdots,q,\ t\in\bb{R}$ y $\Vert z\Vert\leq r$ and 
\[
\begin{aligned}
\bigg\Vert \frac{\partial^{2}f}{\partial x_{i}\partial x_{j}}(t,\theta z(t)+\xi(t),\theta z([t])+\xi([t])) -\frac{\partial^{2}f}{\partial x_{i}\partial x_{j}} & (t,\theta\tilde{z}(t)+\xi(t),\theta\tilde{z}([t])+\xi([t]))\bigg\Vert \\
&\leq~{}2\theta N_{2}(r)\left\Vert z(t)-\tilde{z}(t)\right\Vert
\\
\bigg\Vert \frac{\partial^{2}f}{\partial y_{i}\partial y_{j}}(t,\theta z(t)+\xi(t),\theta z([t])+\xi([t]))-\frac{\partial^{2}f}{\partial y_{i}\partial y_{j}}&(t,\theta\tilde{z}(t)+\xi(t),\theta\tilde{z}([t])+\xi([t]))\bigg\Vert \\
&\leq ~{} 2\theta N_{2}(r)\big\Vert z(t)-\tilde{z}(t)\big\Vert
\\
\bigg\Vert\frac{\partial^{2}f}{\partial x_{i}\partial y_{j}}(t,\theta z(t)+\xi(t),\theta z([t])+\xi([t]))-\frac{\partial^{2}f}{\partial x_{i}\partial y_{j}}&(t,\theta\tilde{z}(t)+\xi(t),\theta\tilde{z}([t])+\xi([t]))\bigg\Vert 
\\
&\leq~{} 2\theta N_{2}(r)\left\Vert z(t)-\tilde{z}(t)\right\Vert
\end{aligned}
\]
where $i,j=1,2,\cdots,q,\ t\in\bb{R}$ and $\Vert z\Vert\leq r$, $N_2(r)$ is bounded and $N_1(r),M_0(\nu)$ can be chosen as non decreasing functions of $r$ and $\nu$ respectively. It follows
\begin{eqnarray*}
\Vert F(t,z(t),z([t]))\Vert \leq 3r^{2}N_1(r) \mbox{ for } \Vert z\Vert\leq r
\end{eqnarray*}
and
\begin{eqnarray*}
\Vert F(t,z(t),z([t]))-F(t,\tilde{z}(t),\tilde{z}([t]))\Vert \leq 6(qrN_1(r)+r^{2}N_2(r))\Vert z-\tilde{z}\Vert
\end{eqnarray*}
for $\Vert z\Vert, \Vert\tilde{z} \Vert\leq r$ and $t\in \bb{R}$.

Then, since the hypotheses of Theorem \ref{teo_perturbada_depca} are satisfied, the theorem follows.

\end{proof}

\end{document}
